% Generated draft: requires mathematical and source review.
\documentclass[11pt]{article}
\usepackage[margin=25mm]{geometry}
\usepackage{fontspec}
\setmainfont{Latin Modern Roman}
\usepackage{amsmath,amssymb,mathtools,unicode-math}
\setmathfont{Latin Modern Math}
\usepackage{enumitem,needspace,xurl,hyperref}
\setlength{\parindent}{0pt}
\setlength{\parskip}{5pt}
\setlength{\emergencystretch}{3em}
\newcommand{\N}{\mathbb{N}}
\newcommand{\Z}{\mathbb{Z}}
\newcommand{\Q}{\mathbb{Q}}
\newcommand{\R}{\mathbb{R}}
\newcommand{\C}{\mathbb{C}}
\newcommand{\F}{\mathbb{F}}
\newcommand{\A}{\mathbb{A}}
\newcommand{\PP}{\mathbb{P}}
\newcommand{\E}{\mathbb{E}}
\newcommand{\eps}{\varepsilon}
\newcommand{\dd}{\,\mathrm{d}}
\begin{document}
\section*{UnsolvedMath ID: 30004438}
\textbf{Source problem code:} catalogue entry
\subsection{Definitions and mathematical statement}
Fix an integer \(d\geq 2\). Let \(\operatorname{Rat}_d(\R)\) denote the space of degree-\(d\) morphisms \(\phi:\PP^1_{\R}\to\PP^1_{\R}\). Concretely, its elements are pairs \([F:G]\) of homogeneous real polynomials of degree \(d\) with no common zero in \(\PP^1(\C)\), modulo multiplication of both polynomials by the same nonzero real scalar. Equip this open subset of \(\PP^{2d+1}(\R)\) with its usual real topology; its real manifold dimension is \(2d+1\). The parameter space here consists of maps themselves, without quotienting by conjugacy.

For \(\phi\in\operatorname{Rat}_d(\R)\), write \(\phi_{\C}\) for its extension to \(\PP^1(\C)\). A point \(z\in\PP^1(\C)\) is periodic if \(\phi_{\C}^{\circ n}(z)=z\) for some integer \(n\geq1\). Define
\[
\mathcal R_d=\left\{\phi\in\operatorname{Rat}_d(\R):\ \text{for every }n\geq1\text{ and every }z\in\PP^1(\C),\quad \phi_{\C}^{\circ n}(z)=z\ \Longrightarrow\ z\in\PP^1(\R)\right\}.
\]
Thus periodic points of every period are included, and the point at infinity is included in both projective lines.

\textbf{Remaining question.} Is it true that, for every integer \(d\geq2\),
\[
\operatorname{int}_{\operatorname{Rat}_d(\R)}(\mathcal R_d)\neq\varnothing?
\]
Equivalently, for each such \(d\), does there exist a nonempty open subset \(U_d\subseteq\operatorname{Rat}_d(\R)\) such that every periodic point of every map in \(U_d\) is real?

The original formulation also asks whether \(\mathcal R_d\) has full dimension, and presents this as equivalent to nonempty interior. The supplied record does not specify the notion of dimension or justify this equivalence. Consequently, the interior question above is precise, while the claimed equivalence with the dimension formulation remains a source qualification requiring review.

The original question additionally asked whether \(\mathcal R_d\neq\varnothing\) for each \(d\geq2\). The supplied curation reports an affirmative answer via Chebyshev maps and labels the overall problem partially solved. It reports that the interior/full-dimension question remains open and that only lower-dimensional families are known. These are supplied literature assessments, not independent conclusions of this draft.
\subsection{Short English statement}
For every degree \(d\geq2\), ask whether there is a nonempty open family of real rational maps for which every complex periodic point lies on the real projective line. The supplied record attributes existence of individual examples to Chebyshev maps. Its asserted equivalence between full dimension and nonempty interior needs clarification.
\subsection{Sources}
\begin{itemize}
\item Local UnsolvedMath registry: \texttt{data/unsolvedmath/problems.json}, supplied record with permanent ID \(30004438\), source problem code \texttt{catalogue entry}, entitled \emph{Real Rational Maps with Exclusively Real Periodic Points}. Collection: Oberwolfach Reports Open Problems; proposed year supplied as \(2020\). No individual proposer is supplied.
\item Supplied report citation: \emph{Real Algebraic Geometry with a View Toward Hyperbolic Programming and Free Probability} (2021), Oberwolfach Reports, pp.\ 639--712, DOI \texttt{10.4171/owr/2020/12}. \url{https://doi.org/10.4171/owr/2020/12}. The background describes Oberwolfach Reports as records of mathematical workshops and their open-problem discussions.
\item Supplied literature source: Kozhasov--Kummer (2020), \url{https://arxiv.org/abs/2004.10003}.
\item Supplied literature source: Oberwolfach Report 12/2020, \url{https://doi.org/10.4171/owr/2020/12}.
\item Literature assessment dated exactly \texttt{2026-08-22}: supplied dated curation, not independently verified here. It reports that existence is settled affirmatively by the Chebyshev morphism, that the nonempty-interior/full-dimension question remains open, and that only lower-dimensional families are known. The background explicitly requests independent verification of time-sensitive status and cited claims.
\item The registry labels the statement \texttt{corrected\_verified} and reports that the official report was checked and the solved existence clause distinguished from the remaining question. These are supplied verification claims; this draft does not claim to have visited or checked the report.
\end{itemize}
\end{document}
